\section{Separation and comparison tools}
\label{comp:tools}

This appendix gives the exact statements behind the numerical comparisons.
The statements concern exact arithmetic. The archived implementations use
floating point, so their reported values need the qualifications in
Section~\ref{comp:numerics}.

\subsection{Family separation on the simplex}

Write $\Delta_S=\{v\in\R_+^S:\mathbf1^Tv=1\}$ for a nonempty
$S\subseteq\{1,2,3,4\}$. For $A\in\Sn4$, define
\begin{equation}\label{comp:simplex-min}
 \mu(A)=\min_{v\in\Delta_{\{1,2,3,4\}}}v^TAv,
 \qquad
 K_S=\begin{pmatrix}A_{SS}&\mathbf1\\\mathbf1^T&0\end{pmatrix}.
\end{equation}
The normalization in \eqref{comp:simplex-min} differs from the Euclidean
normalization in Proposition~\ref{family:separation}; both detect exactly
whether $A$ is copositive.

\begin{proposition}[Fifteen-support separation]\label{comp:family-separation}
For each nonempty support $S$, solve
\begin{equation}\label{comp:bordered}
 K_S\begin{pmatrix}v_S\\\nu\end{pmatrix}
   =\begin{pmatrix}0\\1\end{pmatrix}.
\end{equation}
Retain a solution whenever $v_S\geq0$, and extend it by zero outside $S$.
The least objective value among these feasible stationary solutions is
$\mu(A)$. If a system is singular but has feasible solutions, all of them
have the same objective value. Moreover, a global minimizer occurs on a
support for which $K_S$ is nonsingular. Thus singular systems may be omitted
in exact arithmetic without losing the minimum.
\end{proposition}

\begin{proof}
Every retained vector is feasible for \eqref{comp:simplex-min}, so its
objective is at least $\mu(A)$. A global minimizer exists by compactness.
Choose one, $v^*$, with the smallest support $S$. It lies in the relative
interior of $\Delta_S$, so first-order stationarity on its affine hull gives
$A_{SS}v^*_S=\lambda\mathbf1$. Hence $(v^*_S,-\lambda)$ solves
\eqref{comp:bordered}.

Suppose $K_S$ were singular. There would be a nonzero pair $(d,\eta)$ with
$A_{SS}d+\eta\mathbf1=0$ and $\mathbf1^Td=0$. Necessarily $d\ne0$.
For every $t\in\R$,
\[
 (v^*_S+td)^TA_{SS}(v^*_S+td)
 = (v^*_S)^TA_{SS}v^*_S
   +2t\lambda\mathbf1^Td-t^2\eta\mathbf1^Td
 =\mu(A).
\]
Because $d$ has both positive and negative coordinates and $v^*_S>0$,
one can move in this direction until a coordinate first becomes zero.
The resulting point is feasible and has a smaller support, a contradiction.
Thus $K_S$ is nonsingular and $v^*$ is among the retained candidates.

Finally, if $u,v\in\Delta_S$ solve the stationary equations with
$A_{SS}u=\lambda\mathbf1$ and $A_{SS}v=\rho\mathbf1$, symmetry gives
$\lambda=v^TA_{SS}u=u^TA_{SS}v=\rho$. Their objective values are
therefore equal. This also handles degenerate faces with infinitely many
stationary solutions. No positive-definiteness assumption on $A_{SS}$
is needed: nonminimizing stationary points remain feasible candidates and
cannot lower the minimum below $\mu(A)$.
\end{proof}

For the family of Section~\ref{family:section}, take
$A=B-bb^T$ from \eqref{family:lmi}. Minimization over the unrestricted
parameter $h$ gives $h=-b^Tv$ and family value $v^TAv$.
A negative value in \eqref{comp:simplex-min} therefore gives a violated
family inequality, with $(d_1,d_2,d_3,k)=v$. Any vector $v\geq0$
gives a valid inequality, even if it is not a minimizer. The archived
implementation clips tiny negative coordinates, renormalizes the vector,
and evaluates the resulting feasible candidate. This preserves the
mathematical validity of its parameter choice, but skipping a nearly
singular system can miss a violation. A floating-point value of $\mu(A)$
is consequently a separation diagnostic, not a certificate of copositivity.

\subsection{The five-tetrahedron lift and normalized hull depth}

For $u\in C_3$, put $\bar u=(1,u)^T$, and write
\[
 \mathcal M_3
   =\conv\{\bar u\bar u^T:u\in C_3\}.
\]
This is the matrix form of $\mathcal Q_3$. Let
\begin{equation}\label{comp:uniform}
 M_c=\begin{pmatrix}
 1&1/2&1/2&1/2\\
 1/2&1/3&1/4&1/4\\
 1/2&1/4&1/3&1/4\\
 1/2&1/4&1/4&1/3
 \end{pmatrix},
\end{equation}
the moment matrix of the uniform distribution on $C_3$.
Use the following five tetrahedra; strings denote cube vertices:
\begin{equation}\label{comp:tetrahedra}
 \begin{split}
 &(000,110,101,011),\qquad(100,000,110,101),\\
 &(010,000,110,011),\qquad(001,000,101,011),\\
 &(111,110,101,011).
 \end{split}
\end{equation}
Let $\bar A_p$ have the four homogenized vertices of tetrahedron $p$
as columns. To see that these tetrahedra cover the cube, the central
tetrahedron has barycentric coordinates
\[
 1-\tfrac{x+y+z}{2},\quad\tfrac{x+y-z}{2},\quad
 \tfrac{x+z-y}{2},\quad\tfrac{y+z-x}{2}.
\]
Outside it, one of $x>y+z$, $y>x+z$, $z>x+y$, or $x+y+z>2$
holds. These regions are the four remaining tetrahedra, and their interiors
are disjoint.

The exact lift of \citet{AnstreicherBurer2010} specializes to
\begin{equation}\label{comp:exact-lift}
 \mathcal M_3
 =\left\{\sum_{p=1}^5\bar A_pW_p\bar A_p^T:
      W_p\in\DNN_4,\quad
      \sum_{p=1}^5\mathbf1^TW_p\mathbf1=1\right\}.
\end{equation}
For completeness, $\DNN_4=\CP_4$: factor each $W_p$ into a sum of
$ww^T$ with $w\geq0$. For $w\ne0$, let
$\alpha=\mathbf1^Tw$ and $u$ be the point of tetrahedron $p$ with
$\bar u=\bar A_pw/\alpha$. Then
$\bar A_pww^T\bar A_p^T=\alpha^2\bar u\bar u^T$, and the
normalization in \eqref{comp:exact-lift} says that these weights sum to one.
Conversely, every cube point belongs to a tetrahedron and has nonnegative
barycentric coordinates, yielding a feasible rank-one $W_p$.

\begin{proposition}[Hull depth and exact cut correction]\label{comp:hull-depth}
For $M\in\Sn4$ with $M_{00}=1$, define
\begin{equation}\label{comp:depth}
 \begin{split}
 \delta(M)=\min\quad&\inner{C}{M}\\
 \text{subject to}\quad&\inner{C}{M_c}=1,\\
 &\bar A_p^TC\bar A_p-N_p\succeq0\quad(p=1,\ldots,5),\\
 &N_p\geq0\ \text{entrywise},\qquad\diag(N_p)=0.
 \end{split}
\end{equation}
The minimum is finite and attained. The projected feasible matrices $C$
are exactly the quadratics $q_C(u)=\bar u^TC\bar u$ that are
nonnegative on $C_3$ and have uniform mean one. Consequently,
\[
 \delta(M)<0\quad\Longleftrightarrow\quad
 M\notin\mathcal M_3.
\]
If $\varepsilon\geq\max\{0,-\delta(M)\}$, then
\begin{equation}\label{comp:depth-mixing}
 \frac{M+\varepsilon M_c}{1+\varepsilon}
       \in\mathcal M_3.
\end{equation}
For any symmetric $C$, the shifted quadratic $q_C+s$ is nonnegative
on $C_3$ whenever
\begin{equation}\label{comp:shift}
 s\geq\max\left\{0,-\min_{1\leq p\leq5}
          \mu(\bar A_p^TC\bar A_p)\right\}.
\end{equation}
\end{proposition}

\begin{proof}
On tetrahedron $p$, $q_C(u)=\lambda^T\bar A_p^TC\bar A_p\lambda$
for $\lambda\geq0$ with $\mathbf1^T\lambda=1$. Thus $q_C$ is
nonnegative on the cube exactly when all five matrices
$\bar A_p^TC\bar A_p$ are copositive. The equality
$\COP_4=\mathbb S^4_++(\R_+^{4\times4}\cap\Sn4)$ gives precisely the constraints
in \eqref{comp:depth}; any nonnegative diagonal component can be moved
into the positive semidefinite component.

The cone of cube-nonnegative quadratics is closed. A nonzero element has
strictly positive uniform mean: a continuous nonnegative polynomial with
zero integral vanishes on the cube and hence is the zero polynomial.
Its mean-one slice is compact. Indeed, an unbounded sequence in this
slice, divided by its coefficient norm, would have a subsequence converging
to a nonzero nonnegative quadratic of mean zero, a contradiction.
The slice is nonempty because it contains the constant one. Therefore
\eqref{comp:depth} has a finite attained minimum in $C$, and a
copositive decomposition supplies the auxiliary $N_p$.

Every $M\in\mathcal M_3$ evaluates all nonnegative
quadratics nonnegatively. If $M$ is outside this compact convex hull,
strict separation gives a matrix $D$ and a scalar $a$ with
$\inner{D}{M}<a\leq\inner{D}{\bar u\bar u^T}$ for every cube
point. Subtracting $a$ from the constant term yields a nonzero
cube-nonnegative quadratic with negative value at $M$, because
$M_{00}=1$. Its positive uniform mean permits normalization in
\eqref{comp:depth}. This proves the equivalence.

For every normalized nonnegative quadratic,
$\inner{C}{M+\varepsilon M_c}\geq\delta(M)+\varepsilon\geq0$.
The preceding separation characterization and the leading entry
$1+\varepsilon$ give \eqref{comp:depth-mixing}. Finally, on tetrahedron
$p$, $q_C(u)\geq\mu(\bar A_p^TC\bar A_p)$, proving
\eqref{comp:shift}.
\end{proof}

Formula~\eqref{comp:shift} is an exact certification rule if its five
simplex minima are evaluated exactly, or if certified lower bounds replace
them. Applying the floating-point routine of
Proposition~\ref{comp:family-separation} with a small rounding margin
does not establish this rule in exact arithmetic.

\subsection{An upper bound on triple-local improvement}

Let $y$ be a vector of first and second moments on a fixed pattern, with
every diagonal recorded. For a recorded triple $T$, write $M_T(y)$ for
its normalized $4\times4$ moment matrix. Let $y_c$ be the restriction of
the uniform cube moments to this pattern. Thus each $M_T(y_c)$ equals
\eqref{comp:uniform}.

\begin{proposition}[Triple-local gain bound]\label{comp:gain-bound}
Let $R$ be a convex set of moment vectors containing $y_c$, let $y^*\in R$,
and let $f$ be an affine objective. Let $\mathcal T$ be a set of recorded
triples and assume
\[
 \varepsilon\geq\max_{T\in\mathcal T}
                      \max\{0,-\delta(M_T(y^*))\}.
\]
For each $T$, add any system involving its moments and auxiliaries used
only by that triple, whose projection contains
$\mathcal M_3$. Let $R'$ be the resulting moment feasible set.
Then
\begin{equation}\label{comp:gain}
 \inf_{y\in R'}f(y)
 \leq \frac{f(y^*)+\varepsilon f(y_c)}{1+\varepsilon}
 =f(y^*)+\frac{\varepsilon}{1+\varepsilon}
                         \bigl(f(y_c)-f(y^*)\bigr).
\end{equation}
If $y^*$ minimizes $f$ over $R$, the last term bounds the improvement
of the relaxation value.
\end{proposition}

\begin{proof}
Convexity gives $y_\varepsilon=(y^*+\varepsilon y_c)/(1+\varepsilon)\in R$.
By Proposition~\ref{comp:hull-depth}, $M_T(y_\varepsilon)$ belongs to
$\mathcal M_3$ for every $T\in\mathcal T$. Each added system
has feasible auxiliaries at these moments. Because the auxiliaries are
independent between triples, they can all be chosen at once. Hence
$y_\varepsilon\in R'$. Evaluating $f$ there proves \eqref{comp:gain}.
\end{proof}

This argument applies unchanged to clique moment patterns. It covers the
family blocks, tetrahedral lifts, and the Anstreicher--Puges systems with
their own trilinear moment for each triple. It also covers the disjoint-support
systems when their additional moments are not identified across triples.
It does not cover a common auxiliary such as $x_i^2x_j$ used by two
different triples: local representing measures need not give the same
value to that shared auxiliary.

For the objective $f(m,Y)=\inner{H}{Y}+g^Tm$,
\[
 f(y_c)=\tfrac13\sum_iH_{ii}
        +\tfrac14\sum_{i\ne j}H_{ij}+\tfrac12\sum_i g_i.
\]
An application to a numerical $y^*$ requires actual feasibility for $R$
and a one-sided upper bound on every $-\delta(M_T(y^*))$.
If $L_R$ is a reported lower bound on the baseline value, the comparison
with $L_R$ must also include the margin $f(y^*)-L_R$.

\subsection{Dual residual correction}

\begin{proposition}\label{comp:dual-bound}
Consider $\inf\{c^Tu:Au+s=b,\ s\in K\}$, where every feasible $u$
satisfies $\ell_i\leq u_i\leq h_i$. For any $z\in K^*$, put
$r=c+A^Tz$. Then every feasible point satisfies
\begin{equation}\label{comp:dual-correction}
 c^Tu\geq-b^Tz+\sum_i\min\{r_i\ell_i,r_ih_i\}.
\end{equation}
\end{proposition}

\begin{proof}
Since $\inner{z}{s}\geq0$,
$c^Tu=r^Tu-z^TAu=r^Tu-b^Tz+\inner{z}{s}
\geq r^Tu-b^Tz$. Minimize each component of $r^Tu$ over its interval.
\end{proof}

The implementation projects the approximate dual vector onto the relevant
cones and subtracts rounding allowances before using
\eqref{comp:dual-correction}. The variables have explicit finite bounds:
recorded moments lie in $[0,1]$, the six family auxiliaries lie in
$[0,3]$, and tetrahedral weights lie in $[0,1]$. The family bound is
implied by the block at a baseline-feasible point: each entry of $B$
and $b$ has absolute value at most one, while
$S=B-N-bb^T\succeq0$ has diagonal at most one and hence entries
of absolute value at most one. Thus $N_{ij}=B_{ij}-b_ib_j-S_{ij}\leq3$.
For the tetrahedral lift, nonnegativity and
$\sum_p\mathbf1^TW_p\mathbf1=1$ imply each entry of $W_p$ is at
most one. These restrictions preserve the intended lifts. The reported corrected values are careful
floating-point estimates. They become rigorously certified lower bounds
only if cone membership, residuals, and the rounding allowances are
themselves verified with rigorous arithmetic.
